\chapter{Starting Data for Local Quantum Field Theory}
\label{ch:starting-data-local-qft}

\section{Spacetime And Operator Conventions}

Throughout the opening volume, \(D\) denotes the spacetime dimension. The
symbol \(\mathbb M^D\) denotes \(D\)-dimensional Minkowski affine space. After
a choice of inertial coordinates, \(\mathbb M^D\) is identified with
\(\mathbb R^D\). The metric tensor has components
\[
  \eta_{\mu\nu}=\operatorname{diag}(-1,1,\ldots,1),
\]
and units are chosen so that \(\hbar=c=1\). In the quantum-mechanical
path-integral derivations \(\hbar\) is displayed until the semiclassical
weighting and Wick rotation have been derived; after that point \(\hbar=1\)
is restored unless a dimensional check explicitly reintroduces it.  The
vacuum vector is denoted by \(\vac\); the same printed letter may occur with
subscripts for reconstructed cyclic vectors, but unsubscripted \(\vac\)
always denotes the vacuum of the theory under discussion. A vector
\(x\in\mathbb M^D\) is
spacelike when
\[
  x^2:=\eta_{\mu\nu}x^\mu x^\nu>0.
\]
A test function is a smooth compactly supported complex-valued function on
\(\mathbb M^D\), an element of \(C_c^\infty(\mathbb M^D)\), unless another
test-function space is explicitly declared. A pointlike quantum field
\(\widehat\Phi_A(x)\) is an operator-valued distribution: its primary
operator expression is the smeared field
\[
  \widehat\Phi_A(f)
  =
  \int_{\mathbb M^D}\dd^Dx\,f(x)\widehat\Phi_A(x),
  \qquad
  f\in C_c^\infty(\mathbb M^D).
\]
The integral sign is distributional notation for evaluating the field on the
test function \(f\).  When \(\widehat\Phi_A(f)\) is unbounded, its common dense
domain is part of the field framework being used.  The symbol
\(\mathcal B(\Hilb)\) denotes the algebra of bounded linear operators on a
Hilbert space \(\Hilb\), and \(\mathcal U(\Hilb)\) denotes its unitary group.
The label \(A\) denotes a component in a finite-dimensional representation
space for Lorentz and internal indices, or a composite-operator label, as
specified in the chapter where the field is introduced.

Euclidean correlation distributions are called Schwinger functions when they
are being used as reconstruction data and are then denoted \(S_n\).  In
path-integral perturbation chapters the notation \(G_{E,n}\) denotes the same
kind of Euclidean \(n\)-point distribution before the OS reconstruction
hypotheses have been imposed; \(G_E\) without an \(n\)-index is reserved for a
two-point kernel or its Fourier transform.

The connected proper orthochronous Lorentz group is denoted
\(SO^+(1,D-1)\), and the connected Poincare group is
\(\mathbb R^D\rtimes SO^+(1,D-1)\).  Two subsets
\(\mathcal O_1,\mathcal O_2\subset\mathbb M^D\) are spacelike separated when
\((x-y)^2>0\) for every \(x\in\mathcal O_1\) and
\(y\in\mathcal O_2\).

\section{Local Continuum Theories And Effective Presentations}

The organizing structure of quantum field theory is quantum dynamics together
with spacetime localization.  The word ``localization'' has a precise content:
observables, or field coordinates for observables, are attached to spacetime
regions or test functions; spacetime symmetries move these localized objects
covariantly; and observables localized in spacelike separated regions obey the
appropriate commutativity law.  The exact algebraic, field-theoretic, or
Euclidean form of this structure must always be specified.

Two different mathematical statuses will appear throughout the monograph.
They share locality, but they are not the same assertion.

\begin{definition}[Continuum local quantum field theory]
  A continuum local quantum field theory is a local quantum theory whose
  declared correlation functions, local algebras, or Hilbert-space operators
  exist after the ultraviolet regulator has been removed, in a stated
  topology and within a stated spacetime and state-sector framework.  In a
  Minkowski vacuum presentation this includes Poincare symmetry, a spectrum
  condition, local observables or local field distributions, and
  microcausality.  The phrase ``infinitely many degrees of freedom'' will be
  used in the operational sense that the local data are indexed by
  infinite-dimensional test-function or region data and are not specified by a
  finite list of canonical coordinates.
\end{definition}

\begin{definition}[Effective field theory presentation]
  An effective field theory presentation is a regulated or renormalized local
  prescription for computing observables in a specified range of length,
  energy, and external-state scales.  Its defining data include the fields or
  local variables, the symmetries, the regulator or renormalization
  prescription, the expansion parameters, and the class of observables for
  which the expansion is claimed.  When it is perturbative, it defines an
  asymptotic expansion of Green functions, and of scattering quantities only
  after the nonperturbative scattering framework and LSZ interpretation have
  been supplied.
\end{definition}

Power-counting renormalizability is a property of an effective presentation
near a chosen fixed point or scaling regime: only finitely many local
operators are needed as independent counterterms at a fixed order in the
declared expansion.  It is not, by itself, a proof that the ultraviolet
continuum theory exists.  Power-counting nonrenormalizable presentations are
also local effective theories when the infinitely many allowed operators are
organized by scale and symmetry.

\begin{figure}[H]
  \centering
  \resizebox{0.96\textwidth}{!}{%
  \begin{tikzpicture}[x=1cm,y=1cm,every node/.style={font=\scriptsize}]
    \tikzset{
      root/.style={draw,rounded corners=4pt,line width=0.85pt,
        fill=black!3,minimum width=5.1cm,minimum height=0.7cm,align=center},
      box/.style={draw,rounded corners=4pt,line width=0.8pt,
        minimum width=5.4cm,minimum height=1.0cm,align=center},
      textbox/.style={draw,rounded corners=4pt,line width=0.55pt,
        align=left,inner sep=5pt,text width=5.8cm},
      arrow/.style={-{Latex[length=2mm]},line width=0.8pt}
    }
    \node[root] (root) at (0,2.35)
      {quantum dynamics with spacetime localization};
    \node[box,fill=blue!5] (local) at (-4.0,0.82)
      {continuum local QFT\\data exist after removing the ultraviolet regulator};
    \node[box,fill=orange!7] (eft) at (4.0,0.82)
      {effective field theory presentation\\local expansion valid in a declared scale window};
    \draw[arrow] (root.south west) -- (local.north);
    \draw[arrow] (root.south east) -- (eft.north);

    \node[textbox,anchor=north,fill=blue!2] at (-4.0,-0.15)
      {\(\Hilb\), Poincare representation, spectrum condition\\
       local observables or local field distributions\\
       microcausality for spacelike separated supports\\[2pt]
       guiding examples: low-dimensional constructive scalar theories;
       scaling limits of lattice models near criticality; Yang--Mills/QCD as
       physical continuum local gauge theories};

    \node[textbox,anchor=north,fill=orange!3] at (4.0,-0.15)
      {regulated path integral or equivalent local prescription\\
       symmetry-controlled local operator expansion\\
       long-distance or low-energy observables in the stated regime\\[2pt]
       examples: \(\phi^4\) theory in \(D=4\), QED, and the Standard Model as
       power-counting renormalizable perturbative presentations; chiral
       perturbation theory and perturbative quantum gravity as
       nonrenormalizable effective presentations};
  \end{tikzpicture}
  }
  \caption{The opening distinction between continuum local QFT data and
  effective field theory presentations.  Both sides use locality; the
  difference is the mathematical status of the construction and the range of
  claims made from it.}
  \label{fig:opening-local-qft-eft}
\end{figure}

\section{Opening Framework}

The first construction takes place in a flat-spacetime vacuum sector. The
following definition gives the data used in the opening volume.

\begin{definition}[Vacuum Minkowski local quantum theory]
\label{def:vacuum-minkowski-local-qft}
In the opening framework, a vacuum Minkowski local quantum theory consists of:
\begin{enumerate}
  \item a complex Hilbert space \(\Hilb\);
  \item a spacetime symmetry convention.  Let \(G_{\rm Lor}\) be
  \(SO^+(1,D-1)\), or the appropriate double cover when spinor fields are
  included, and let
  \(\pi:G_{\rm Lor}\to SO^+(1,D-1)\) be the covering projection.  The connected
  Poincare group used here is
  \[
    \mathcal P=\mathbb R^D\rtimes G_{\rm Lor},
    \qquad
    (a,\Lambda)(b,\Gamma)
    =
    (a+\pi(\Lambda)b,\Lambda\Gamma).
  \]
  It acts on \(\mathbb M^D\) from the left by
  \[
    g x=\pi(\Lambda)x+a,
    \qquad g=(a,\Lambda)\in\mathcal P,
  \]
  and \(g\mathcal O:=\{gx:x\in\mathcal O\}\) for a region
  \(\mathcal O\subset\mathbb M^D\);
  \item a strongly continuous unitary representation
  \[
    U:\mathcal P\to\mathcal U(\Hilb),
    \qquad
    g\mapsto U(g),
  \]
  as a left representation, \(U(g_1g_2)=U(g_1)U(g_2)\), where
  \(\mathcal U(\Hilb)\) is the unitary group of \(\Hilb\);
  \item self-adjoint translation generators \(P_\mu\), defined by
  \(U(a,1)=\exp(\ii a^\mu P_\mu)\).  Let \(E_P\) be the joint
  projection-valued spectral measure of \(P=(P_0,\ldots,P_{D-1})\).  The
  spectrum condition is
  \[
    \operatorname{supp}E_P\subset
    \overline V_+
    =
    \{p\in\mathbb R^D:p^0\ge0,\ \eta_{\mu\nu}p^\mu p^\nu\le0\};
  \]
  \item a Poincare-invariant unit vector \(\vac\in\Hilb\),
  \[
    U(g)\vac=\vac,\qquad g\in\mathcal P.
  \]
  Vacuum uniqueness in the chosen sector means that the zero-momentum spectral
  projection is one-dimensional:
  \[
    E_P(\{0\})\Hilb
    =
    \mathbb C\vac;
  \]
  \item an assignment
  \[
    \mathcal O\longmapsto\Obs(\mathcal O)\subset\mathcal B(\Hilb)
  \]
  from bounded open regions \(\mathcal O\subset\mathbb M^D\) to unital
  concrete \(*\)-subalgebras of the bounded operators on \(\Hilb\), satisfying
  covariance, isotony, and locality:
  \[
    U(g)^{-1}\Obs(\mathcal O)U(g)
    =
    \Obs(g\mathcal O),
    \qquad g\in\mathcal P,
  \]
  \[
    \mathcal O_1\subset\mathcal O_2
    \quad\Longrightarrow\quad
    \Obs(\mathcal O_1)\subset\Obs(\mathcal O_2),
  \]
  and
  \[
    [A,B]=0
    \quad
    \text{for }A\in\Obs(\mathcal O_1),\ B\in\Obs(\mathcal O_2)
  \]
  whenever \(\mathcal O_1\) and \(\mathcal O_2\) are spacelike separated, with
  the graded commutator
  \[
    [A,B]_{\rm gr}=AB-(-1)^{|A||B|}BA
  \]
  used for homogeneous odd operators after a fermion parity grading and the
  parities \(|A|,|B|\in\mathbb Z/2\mathbb Z\) have been specified;
  \item vacuum cyclicity for the bounded local observable content.  Let
  \(\Obs_{\rm loc}^{\rm alg}\) denote the unital \(*\)-algebra generated in
  \(\mathcal B(\Hilb)\) by \(\bigcup_{\mathcal O}\Obs(\mathcal O)\), where the
  union is over bounded open regions.  Then
  \[
    \overline{\Obs_{\rm loc}^{\rm alg}\vac}^{\,\|\cdot\|_{\Hilb}}
    =
    \Hilb;
  \]
  \item when unbounded field coordinates are part of the presentation, a
  specified field-domain datum
  \[
    (\mathcal D_\Phi,\{\Phi_A\}_{A\in I}).
  \]
  Here \(I\) is the field-label set, \(\mathcal D_\Phi\subset\Hilb\) is a dense
  linear subspace with \(\vac\in\mathcal D_\Phi\), \(U(g)\mathcal D_\Phi\subset
  \mathcal D_\Phi\) for all \(g\in\mathcal P\), and each smeared field is a
  linear operator
  \[
    \Phi_A(f):\mathcal D_\Phi\to\mathcal D_\Phi,
    \qquad f\in\mathcal T_A,
  \]
  where \(\mathcal T_A\) is the declared test-function space for that field
  component.  Finite products of such smeared fields are defined on
  \(\mathcal D_\Phi\).  If adjoints are used, then
  \(\mathcal D_\Phi\subset\operatorname{Dom}(\Phi_A(f)^*)\) and the adjoint
  fields preserve \(\mathcal D_\Phi\).  For every
  \(\psi,\chi\in\mathcal D_\Phi\), the matrix elements of finite smeared
  products are continuous multilinear functionals of their test functions.
  Field covariance, field locality, and any relation between
  \(\Phi_A(f)\) with \(\operatorname{supp}f\subset\mathcal O\) and
  \(\Obs(\mathcal O)\) are imposed on this same domain or by a specified
  affiliation to the chosen local observable algebra.
\end{enumerate}
\end{definition}

The covariance convention has the following infinitesimal form.  With
\(U(a,1)=\exp(\ii a^\mu P_\mu)\), a coordinate-labelled scalar field obeys
\[
  U(a,1)^{-1}\widehat\Phi(x)U(a,1)
  =
  \widehat\Phi(x+a),
  \qquad
  [P_\mu,\widehat\Phi(x)]
  =
  \ii\,\partial_\mu\widehat\Phi(x)
\]
as an identity of distributions on a common invariant domain.  Equivalently,
conjugation by \(U(a,1)\) shifts the coordinate label by \(-a\).  The local-net
covariance law above uses the same passive coordinate convention: conjugating
an observable by \(U(g)^{-1}\), with \(g=(a,\Lambda)\), sends its localization
region to \(g\mathcal O=\pi(\Lambda)\mathcal O+a\).

In this opening convention \(\Obs(\mathcal O)\) is a concrete algebra of
bounded observables on the chosen Hilbert space and is not assumed to be norm
closed or weakly closed unless this is stated.  If a \(C^*\)-algebra or von
Neumann algebra is needed, the relevant closure or bicommutant is included in
the datum or constructed explicitly.  This convention is sufficient for the
first dependency layer: localization, covariance, the spectrum condition,
vacuum matrix elements, and comparison with field coordinates.

This framework describes a vacuum representation on flat spacetime. Thermal
states, boundary conditions, curved spacetimes, and non-vacuum superselection
sectors are introduced by changing the stated data: the state space, spacetime
category, covariance group, region assignment, or representation sector.
Existing Wightman, Euclidean, algebraic, constructive, perturbative, and
factorization frameworks will be introduced as precise frameworks with their
own domains, beginning with the Wightman and algebraic formulations.

\section{Framework Presentations And Comparison Maps}

The opening data admit several mathematically precise presentations.  Each
presentation emphasizes a different class of objects, and each comparison map
carries hypotheses.

\begin{definition}[Wightman field presentation]
  A Wightman field presentation consists of Hilbert-space data
  \[
    (\Hilb,\vac,U,\mathcal D,\{\Phi_a\}_{a\in A})
  \]
  where \(A\) is an index set for fields, \(\mathcal D\subset\Hilb\) is a
  dense common invariant domain containing \(\vac\), and each \(\Phi_a\) is an
  operator-valued distribution on a declared invariant test-function space
  \(\mathcal T\), with common domain \(\mathcal D\).  In the standard
  Wightman case \(\mathcal T=\mathcal S(\mathbb M^D)\), so the vacuum
  distributions below are tempered.  The domain assumptions are
  \[
    \vac\in\mathcal D,\qquad
    U(g)\mathcal D\subset\mathcal D\quad(g\in\mathcal P),\qquad
    \Phi_a(f)\mathcal D\subset\mathcal D\quad(a\in A,\ f\in\mathcal T),
  \]
  and, for every \(n\ge1\),
  \[
    \Phi_{a_1}(f_1)\cdots\Phi_{a_n}(f_n)\mathcal D
    \subset
    \mathcal D,
    \qquad a_i\in A,\quad f_i\in\mathcal T.
  \]
  Matrix elements of finite products are required to be continuous multilinear
  functionals of the test functions:
  \[
    (f_1,\ldots,f_n)
    \longmapsto
    \left\langle\psi,
    \Phi_{a_1}(f_1)\cdots\Phi_{a_n}(f_n)\chi
    \right\rangle
    \in\mathbb C,
    \qquad \psi,\chi\in\mathcal D .
  \]

  Let \(g=(a,\Lambda)\in\mathcal P\), \(gx=\pi(\Lambda)x+a\), and
  \(f_g(y)=f(g^{-1}y)\).  If the fields carry a finite-dimensional Lorentz
  representation \(R\) of \(G_{\rm Lor}\), covariance is the smeared identity
  \[
    U(g)^{-1}\Phi_a(f)U(g)
    =
    R(\Lambda)_a{}^b\,\Phi_b(f_g),
  \]
  with summation over repeated component labels.  In point-field notation this
  is
  \[
    U(g)^{-1}\Phi_a(x)U(g)=R(\Lambda)_a{}^b\Phi_b(gx).
  \]
  The representation \(U\) satisfies the spectrum condition of
  Definition~\ref{def:vacuum-minkowski-local-qft}: if \(E_P\) is the joint
  spectral measure of the translation generators, then
  \[
    \operatorname{supp}E_P\subset\overline V_+.
  \]

  Locality is imposed after smearing:
  \[
    [\Phi_a(f),\Phi_b(h)]_{\rm gr}=0
  \]
  whenever \(\operatorname{supp}f\) and \(\operatorname{supp}h\) are spacelike
  separated, with the ordinary commutator in the purely bosonic case.

  Its primary correlation data are the vacuum distributions whose distribution
  kernels are written
  \[
    W_{a_1\cdots a_n}(x_1,\ldots,x_n)
    =
    \bra{\vac}
    \Phi_{a_1}(x_1)\cdots\Phi_{a_n}(x_n)
    \ket{\vac}.
  \]
  Equivalently,
  \[
    W_{a_1\cdots a_n}(f_1,\ldots,f_n)
    =
    \left\langle\vac,
    \Phi_{a_1}(f_1)\cdots\Phi_{a_n}(f_n)\vac
    \right\rangle .
  \]
  Translation invariance gives a distribution in difference variables.  With
  \[
    \xi_j=x_{j+1}-x_j,\qquad
    \mathcal W_{a_1\cdots a_n}(\xi_1,\ldots,\xi_{n-1})
    =
    W_{a_1\cdots a_n}
    (0,\xi_1,\xi_1+\xi_2,\ldots,\xi_1+\cdots+\xi_{n-1}),
  \]
  the Wightman spectral condition is the support equation
  \[
    \operatorname{supp}
    \widehat{\mathcal W}_{a_1\cdots a_n}
    \subset
    \overline V_+^{\,n-1}.
  \]
  Positivity of the hierarchy means the following.  For every finite sequence
  of test tensors in the declared test-function space
  \[
    F=\{F^{(n)}_{a_1\cdots a_n}\}_{0\leq n\leq N},
    \qquad
    F^{(0)}\in\mathbb C,
  \]
  define
  \[
    \Psi_F
    =
    F^{(0)}\vac
    +
    \sum_{n=1}^N
    \sum_{a_1,\ldots,a_n}
    \int F^{(n)}_{a_1\cdots a_n}(x_1,\ldots,x_n)
    \Phi_{a_1}(x_1)\cdots\Phi_{a_n}(x_n)\vac\,
    \dd^Dx_1\cdots\dd^Dx_n .
  \]
  The condition is \(\langle\Psi_F,\Psi_F\rangle\geq0\), with the inner product
  evaluated by the distributions \(W_n\):
  \[
    \langle\Psi_F,\Psi_F\rangle\ge0.
  \]
  Vacuum cyclicity is the density equation
  \[
    \overline{
    \operatorname{span}\{
    \Phi_{a_1}(f_1)\cdots\Phi_{a_n}(f_n)\vac
    : n\ge0,\ f_i\in\mathcal T,\ a_i\in A
    \}
    }^{\,\|\cdot\|_{\Hilb}}
    =
    \Hilb .
  \]
\end{definition}

\begin{definition}[Local-net presentation]
  A local-net presentation consists of a specified class
  \(\mathcal K(\mathbb M^D)\) of open spacetime regions, a quasilocal
  \(*\)-algebra or \(C^*\)-algebra \(\Obs_{\rm qloc}\), and an assignment
  \[
    \mathcal O\longmapsto\Obs(\mathcal O)\subset\Obs_{\rm qloc}
  \]
  for \(\mathcal O\in\mathcal K(\mathbb M^D)\).  The assignment satisfies
  isotony, locality, covariance under the declared spacetime symmetry group,
  and any additional structural conditions imposed in the problem under
  study, such as additivity, Haag duality, or the time-slice property.  A
  state is a normalized positive linear functional
  \(\omega:\Obs_{\rm qloc}\to\mathbb C\); it gives a Hilbert-space
  representation by the GNS construction.
\end{definition}

\begin{definition}[Euclidean correlation presentation]
  A Euclidean correlation presentation consists of Schwinger distributions
  \[
    S_{a_1\cdots a_n}(x_1,\ldots,x_n),
    \qquad x_i\in\mathbb R^D,
  \]
  an adjoint involution \(a\mapsto a^\dagger\) on component labels, Euclidean
  covariance, permutation symmetry or graded permutation symmetry, and stated
  regularity and growth assumptions.  Choose Euclidean coordinates
  \(x=(\tau,\mathbf x)\) and let
  \[
    \theta(\tau,\mathbf x)=(-\tau,\mathbf x)
  \]
  be time reflection.  Reflection positivity is the requirement that for every
  finite sequence of test tensors
  \[
    F=\{F^{(n)}_{a_1\cdots a_n}\}_{0\leq n\leq N},
    \qquad
    \operatorname{supp}F^{(n)}
    \subset(\mathbb R_+\times\mathbb R^{D-1})^n,
  \]
  the quadratic form
  \[
  \begin{aligned}
    \langle F,F\rangle_{\rm OS}
    &=
    \sum_{m,n=0}^N
    \sum_{\substack{a_1,\ldots,a_m\\ b_1,\ldots,b_n}}
    \int
    \overline{F^{(m)}_{a_1\cdots a_m}(x_1,\ldots,x_m)}
    F^{(n)}_{b_1\cdots b_n}(y_1,\ldots,y_n)
    \\
    &\qquad\qquad\qquad\quad
    \times
    S_{a_m^\dagger\cdots a_1^\dagger b_1\cdots b_n}
    (\theta x_m,\ldots,\theta x_1,y_1,\ldots,y_n)
    \prod_{i=1}^m\dd^Dx_i\prod_{j=1}^n\dd^Dy_j
  \end{aligned}
  \]
  is nonnegative.  The convention \(S_0=1\) is used in the terms with
  \(m=0\) or \(n=0\).  Under the full Osterwalder--Schrader hypotheses,
  quotienting by the null space of this form and completing reconstructs the
  Lorentzian Hilbert space; analytic continuation then gives Wightman boundary
  values.
\end{definition}

The comparison from fields to a local net is obtained, when the needed
operator-domain and affiliation statements hold, by assigning to a region
\(\mathcal O\) the von Neumann algebra generated by bounded functions,
resolvents, or other bounded functional-calculus representatives of smeared
fields \(\Phi_a(f)\) with \(\operatorname{supp}f\subset\mathcal O\).  The
comparison from a local net to a Hilbert-space theory is the GNS construction
applied to a state, followed by the representation of the local algebras on
the GNS Hilbert space.  The comparison from Euclidean correlations to
Lorentzian fields is OS reconstruction followed by analytic continuation to
Wightman boundary values.

\begin{figure}[H]
  \centering
  \resizebox{0.92\textwidth}{!}{%
  \begin{tikzpicture}[x=1cm,y=1cm,every node/.style={font=\scriptsize}]
    \tikzset{
      box/.style={draw,rounded corners=4pt,line width=0.8pt,fill=blue!4,
        minimum width=3.6cm,minimum height=0.92cm,align=center},
      arrow/.style={-{Latex[length=2mm]},line width=0.85pt},
      note/.style={align=center,font=\scriptsize}
    }
    \node[box] (w) at (-4.4,1.55)
      {Wightman fields\\\(\Phi_a(f),\,W_n\)};
    \node[box] (a) at (4.4,1.55)
      {local nets\\\(\mathcal O\mapsto\Obs(\mathcal O)\)};
    \node[box] (e) at (0,-2.15)
      {Euclidean data\\\(S_n\)};
    \node[box,fill=green!7] (h) at (0,0.05)
      {Hilbert-space vacuum sector\\\((\Hilb,\vac,U)\)};
    \draw[arrow] (w)--(a)
      node[midway,above] {affiliation};
    \draw[arrow] (a.south west) to[out=-155,in=35]
      node[pos=0.56,above,sloped] {GNS} (h.east);
    \draw[arrow] (e)--(h)
      node[midway,right] {OS};
    \draw[arrow] (h.west) to[out=160,in=-30]
      node[pos=0.48,below,sloped] {domains} (w.south east);
    \draw[arrow] (w.south) to[out=-95,in=170]
      node[pos=0.48,left] {analytic continuation} (e.west);
    \node[note] at (0,-3.35)
      {Each arrow includes hypotheses; the comparison data are part of the theory being used.};
  \end{tikzpicture}
  }
  \caption{Three nonperturbative presentations of local QFT and the principal
  comparison maps among them.}
  \label{fig:opening-framework-comparison}
\end{figure}

These definitions make the foundational objects explicit.  A later
calculation may use fields, path integrals, local algebras, or Euclidean
correlators; the framework in which the calculation lives and the comparison
maps needed to interpret it must be stated.

\section{Analytic Status of the Basic Objects}

The opening framework deliberately separates several levels of mathematical
status.  This separation is part of the subject, because the same formula may
denote different objects in different constructions.

\begin{definition}[Regulated theory]
  A regulated theory is a finite-dimensional or cutoff-defined construction
  with a specified state space or configuration space, algebra of variables,
  covariance or symmetry action, and expectation functional
  \(\langle\cdot\rangle_\Lambda\), where \(\Lambda\) denotes the regulator
  data.  Examples include a finite lattice, a finite spatial volume with
  momentum cutoff, or a finite-dimensional time-sliced integral.  All
  integrals and operator products in a regulated theory are ordinary objects
  in the regulator-dependent category.
\end{definition}

\begin{definition}[Continuum correlation theory]
  A continuum correlation theory is a hierarchy of distributions, Schwinger
  functions, Wightman functions, or time-ordered distributions obtained as a
  limit of regulated theories or specified directly by axioms.  The topology
  of convergence is part of the datum: distributional convergence of
  correlators, convergence of moments, weak convergence of measures, or
  another declared topology.
\end{definition}

\begin{definition}[Hilbert-space or algebraic completion]
  A Hilbert-space or algebraic completion is a reconstruction of operators,
  states, local algebras, and symmetries from correlation data or from an
  abstract local net.  Positivity is the bridge: Hilbert spaces arise from
  positive inner products, while local observable algebras arise from
  positive states and their GNS representations.
\end{definition}

Perturbation theory has its own status.  A perturbative statement is an
asymptotic expansion in a specified parameter, with coefficients computed from
regulated distributions and then renormalized.  The existence of the
perturbative coefficients and the existence of a nonperturbative continuum
object are distinct assertions.  A perturbative assertion is specified when
the observable being expanded, the expansion parameter, the renormalization
prescription, and the error class or asymptotic meaning have been stated.

These levels of status will be tracked separately in later chapters.  Spectral
measures, scattering states, renormalization-group flows, anomalies,
conformal data, and gauge-invariant observables are different constructions
related by theorems, controlled approximations, or comparison problems.

\section{Fields As Local Distributional Data}

A field presentation assigns localized distributional operators to test
functions. For a test function \(f\in C_c^\infty(\mathbb M^D)\) with
\(\operatorname{supp}f\subset\mathcal O\),
\[
  \widehat\Phi_A(f)
  =
  \int_{\mathbb M^D}\dd^Dx\,f(x)\widehat\Phi_A(x)
\]
is localized in \(\mathcal O\), in the sense that the smeared operator is
affiliated with the local algebra \(\Obs(\mathcal O)\) or belongs to a common
dense invariant field domain associated with that algebra. The precise
operator-domain statement will be made in the field framework of the next
chapter.

The field presentation is used in two constructions. First, it gives concrete
generators or affiliated operators for examples. Second, it gives correlation
functions:
\[
  W_{A_1\cdots A_n}(x_1,\ldots,x_n)
  =
  \bra{\vac}
  \widehat\Phi_{A_1}(x_1)\cdots
  \widehat\Phi_{A_n}(x_n)
  \ket{\vac},
\]
again as distributions. Under suitable positivity, covariance, spectral, and
domain assumptions, such correlation functions can reconstruct the Hilbert
space and fields. Reconstruction theorems will be used later as comparison
results with stated hypotheses.

\section{Additional Constructions And Hypotheses}

Several familiar objects are constructed from the framework after adding
further hypotheses.

\begin{itemize}
  \item A stable one-particle sector is a spectral subspace of the translation
  generators associated with an isolated positive-energy mass hyperboloid
  \(p^2=-m^2\), \(p^0>0\), possibly with finite internal or spin
  multiplicity.
  \item An S-matrix is constructed from incoming and outgoing asymptotic
  multi-particle states and wave operators comparing asymptotic Hilbert spaces
  with the interacting Hilbert space.  This is an additional long-time
  dynamical construction.
  \item LSZ reduction is a theorem relating S-matrix elements to residues or
  pole parts of time-ordered correlation functions, after the relevant
  asymptotic states and field-particle overlaps have been identified.
  \item A path-integral presentation is a construction or computational
  scheme for correlation functions after a regulator, contour, measure, or
  perturbative prescription has been specified. The symbol
  \(\int D\phi\,\ee^{\ii S[\phi]}\) receives its meaning from that additional
  specification.
\end{itemize}

This dependency order is part of the construction. Perturbative Feynman graphs
may be introduced as an expansion for Green functions before scattering is
discussed. Diagrams for S-matrix elements acquire that interpretation after
asymptotic states and LSZ reduction have been established.

\section{The First Spectral Question}

Translation invariance permits simultaneous spectral analysis of the
generators \(P_\mu\).  Let \(E_P(\Delta)\) denote the joint spectral projection
of \(P=(P_0,\ldots,P_{D-1})\) associated with a Borel set
\(\Delta\subset\mathbb R^D\).  The joint spectrum
\(\operatorname{Spec}(P)\) is the support of this projection-valued measure.
The spectrum condition is
\[
  \operatorname{Spec}(P)\subset
  \overline V_+
  =
  \{p\in\mathbb R^D:p^0\ge0,\ p^2\le0\}.
\]
For a scalar Wightman field in the vacuum Hilbert-space framework, positivity
of the Hilbert-space inner product, translation invariance, and the spectrum
condition imply that the two-point function admits a positive spectral
measure on invariant mass. In the common case where the measure can be
written in terms of a density \(\rho\), the time-ordered two-point function
has the form
\[
  \bra{\vac}T\{\widehat\phi(x)\widehat\phi(0)\}\ket{\vac}
  =
  \int_0^\infty \dd\mu^2\,\rho(\mu^2)\Delta_F(x;\mu^2),
  \qquad
  \rho(\mu^2)\ge0.
\]
Here \(\Delta_F(x;\mu^2)\) denotes the free scalar Feynman distribution of
mass parameter \(\mu\).  If the spectral measure is not absolutely
continuous, the integral is written with respect to the positive measure
\(\dd\rho(\mu^2)\) rather than a density.
This formula will be derived from the stated Hilbert-space assumptions in
Chapter~\ref{ch:kallen-lehmann}.  It is the first result in the text that
relates local two-point data to particle content: isolated atoms of the spectral
measure correspond to stable one-particle states, while continuous spectral
support corresponds to multi-particle and other continuum contributions.

\section{Logical Order Of The Opening Construction}

The opening construction will proceed as follows:
\[
\begin{gathered}
  \text{Hilbert space, symmetry, spectrum, vacuum, locality}
  \\
  \longrightarrow
  \text{local fields and Wightman functions}
  \longrightarrow
  \text{path-integral constructions of Green functions}
  \\
  \longrightarrow
  \text{K\"all\'en--Lehmann spectral representation}
  \longrightarrow
  \text{perturbative Green functions}
  \\
  \longrightarrow
  \text{Haag--Ruelle asymptotic states}
  \longrightarrow
  \text{LSZ}
  \longrightarrow
  \text{perturbative scattering amplitudes}.
\end{gathered}
\]
Each arrow indicates an added construction or theorem. The construction will
use each arrow only after its assumptions have been stated.

The first detailed computations follow three dependency bands.  First, finite
and regulated Lagrangian quantum mechanics supplies path integrals,
regularization, perturbation theory for correlation functions, diagrammatics,
and local counterterms.  Second, relativistic scalar fields supply the first
local field examples, Green functions, spectral representations, asymptotic
states, the nonperturbative scattering operator, and LSZ.  Third, spin and
gauge fields require representation theory, fermionic and gauge-field
quantization, gauge fixing, ghosts, BRST cohomology, and applications such as
QED precision observables.  Perturbative Feynman rules for scattering
amplitudes belong after the second band has supplied their interpretation.
